\subsection{弗雷德霍姆择一定理}

命题 4（弗雷德霍姆择一定理）. --- 设 E 是巴拿赫空间，h 是 E 的一个紧自同态。设 \(\lambda\) 是 K-\{0\} 的一个元素。设 F 是 \(1_{\mathrm{E}} - \lambda h\) 的核，G 是 \(1_{\mathrm{E}'} - \lambda^{t}h\) 的核。

\begin{enumerate}
\def\labelenumi{\alph{enumi})}
\item
  空间 F 与 G 有相同的有限维数 n ≥ 0；
\item
  对 \(y \in \mathrm{E}\) ，存在 \(x \in \mathrm{E}\) 使 \(x - \lambda h(x) = y\) ，当且仅当 \(\langle y, \ell \rangle = 0\) （对一切 \(\ell \in \mathrm{G}\) ）；
\item
  对 \(\ell \in \mathrm{E}'\) ，存在 \(f \in \mathrm{E}'\) 使 \(f - \lambda^t h(f) = \ell\) ，当且仅当 \(\langle x, f \rangle = 0\) （对一切 \(x \in \mathrm{F}\) ）。
\end{enumerate}

特别地，下列条件中恰有一个成立：

\begin{enumerate}
\def\labelenumi{(\roman{enumi})}
\item
  空间 F 非零；
\item
  对每个 \(y \in \mathrm{E}\) ，存在 \(x \in \mathrm{E}\) 使 \(x - \lambda h(x) = y\) 。
\end{enumerate}

通过把 h 换成 \(\lambda h\) ，我们化归到 \(\lambda = 1\) 的情形。

自同态 \(^{t}h\) 由 III, p. 9 的推论 1 是紧的。于是 E 的自同态 \(w = 1_{\mathrm{E}} - h\) 与 \(w' = 1_{\mathrm{E}'} - ^{t}h\) （\(\mathrm{E}'\) 的）都是里斯自同态（III, p. 75 的命题 2 的推论 2）。特别地，它们的核 F 与 G 是有限维的。由于 \(w' = ^{t}w\) 且 w 的指标为零，由 III, p. 43 的公式 (3)，F 的维数等于 G 的维数。

对 E 的元素 \(y\) ，方程 \(x - h(x) = y\) 有解 \(x \in \mathrm{E}\) 当且仅当 \(x\) 属于 \(w\) 的像。由于 \(w\) 的余核的对偶与 \(\operatorname{Ker}(^t w)\) 的核 G 等同（EVT, IV, p. 27, 命题 2），这件事成立当且仅当 \(\langle y, \ell \rangle = 0\) （对一切 \(\ell \in \mathrm{G}\) ）。

对 E' 的元素 \(\ell\) ，方程 \(f - {}^t h(f) = \ell\) 有解 \(f \in \mathrm{E}'\) 当且仅当 \(f\) 属于 \(w'\) 的像。由于 \({}^t w\) 的像是 \(w\) 的核 F 的正交（EVT, IV, p. 27, 命题 2），这件事成立当且仅当 \(\langle x, f \rangle = 0\) （对一切 \(x \in \mathrm{F}\) ）。

例. --- 设 X 是紧空间，\(\mu\) 是 X 上的一个测度（依 K 等于 R 或 C 而为实或复），N: \(X \times X \to K\) 是一个连续函数。

用 E 表示巴拿赫空间 \(\mathcal{C}(\mathrm{X};\mathrm{K})\) 。假设给定了函数 \(a\in\mathrm{E}\) 与 \(\lambda\in\mathrm{K}-\{0\}\) 。我们研究积分方程的解 \(f\in\mathrm{E}\) 的存在性与唯一性问题

\[
f (x) - \lambda \int_ {\mathrm{X}} \mathrm{N} (x, y) f (y) d \mu (y) = a (x) \quad (x \in \mathrm{X}).\tag{3}
\]

为此，引入集 \(F_{\lambda}\) （由相应的齐次方程的解 \(f \in E\) 组成）

\[
f (x) - \lambda \int_ {\mathrm{X}} \mathrm{N} (x, y) f (y) d \mu (y) = 0 \qquad (x \in \mathrm{X}),\tag{4}
\]

以及集 \(G_{\lambda}\) （由 (4) 的转置齐次方程的解 \(g \in E\) 组成），即

\[
g (y) - \lambda \int_ {\mathrm{X}} \mathrm{N} (x, y) g (x) d \mu (x) = 0 \qquad (y \in \mathrm{X}).\tag{5}
\]

若 (3) 的解集非空，则它是 E 的一个以 \(F_{\lambda}\) 为方向的仿射子空间。

对 \(f\in \mathrm{E}\) 与 \(x\in \mathrm{X}\) ，令

\[
h (f) (x) = \int_ {\mathrm{X}} \mathrm{N} (x, y) f (y) d \mu (y);
\]

如此定义的映射 \(f \mapsto h(f)\) 是巴拿赫空间 E 的一个紧自同态（III, p. 26 的命题 2 的推论）。E 的对偶 E' 由 X 上的测度组成（依 K 等于 R 或 C 而为实或复）（INT, III, p. 47, § 1, 第 3 目, 定义 2）。h 的转置 \(^{t}h\) 是 E' 的自同态，由关系式 \(\langle f, ^{t}h(m)\rangle = \langle h(f), m\rangle\) （对 \(f \in E\) 与 \(m \in E'\) ）刻画，即

\[
\begin{array}{r l} \int_ {\mathrm{X}} f (x) d (^ {t} h (m)) (x) & = \int_ {\mathrm{X}} h (f) (x) d m (x) \\ & = \int_ {\mathrm{X}} \Big (\int_ {\mathrm{X}} \mathrm{N} (x, y) f (y) d \mu (y) \Big) d m (x) \\ & = \int_ {\mathrm{X}} \Big (\int_ {\mathrm{X}} \mathrm{N} (x, y) d m (x) \Big) f (y) d \mu (y) \end{array}\tag{6}
\]

（对 \(m \in E'\) 与 \(f \in E\) ；INT, III, p. 84, § 4, 第 1 目, 定理 2）。

引理 4. --- 从 \(G_{\lambda}\) 到 \(E'\) 中的线性映射（由 \(g \mapsto g \cdot \mu\) 定义）是单射的，其像就是 \(1_{E'} - \lambda^{t}h\) 的核。

若 \(g \in G_{\lambda}\) 满足 \(g \cdot \mu = 0\) ，则公式 (5) 蕴含 \(g = 0\) ，故映射 \(g \mapsto g \cdot \mu\) 是单射的。

设 \(m \in \mathrm{E}'\) 是属于 \(1_{\mathrm{E}'} - \lambda^t h\) 的核的一个测度。由关系式 (6)，我们有

\[
\int_ {\mathrm{X}} f (x) d m (x) = \lambda \int_ {\mathrm{X}} \left(\int_ {\mathrm{X}} \mathrm{N} (x, y) d m (x)\right) f (y) d \mu (y)\tag{7}
\]

（对每个连续函数 \(f: \mathrm{X} \to \mathrm{K}\) ）。于是测度 \(m\) 等于测度 \(g \cdot \mu\) ，其中 \(g\) 是由下式定义的从 \(\mathrm{X}\) 到 \(\mathrm{K}\) 中的连续函数

\[
g (y) = \lambda \int_ {\mathrm{X}} \mathrm{N} (x, y) d m (x)\tag{8}
\]

（对一切 \(y \in X\) ）。于是公式 (8) 蕴含函数 g 属于 \(G_{\lambda}\) 。

反之，设 \(g: X \to K\) 是属于 \(G_{\lambda}\) 的一个连续函数，令 \(m = g \cdot \mu\) 。于是对每个连续函数 \(f: X \to K\) ，我们有

\[
\begin{array}{l} \int_ {\mathrm{X}} f (y) d m (y) = \int_ {\mathrm{X}} g (y) f (y) d \mu (y) \\ \qquad = \lambda \int_ {\mathrm{X}} \Big (\int_ {\mathrm{X}} \mathrm{N} (x, y) g (x) d \mu (x) \Big) f (y) d \mu (y) \\ \qquad = \lambda \int_ {\mathrm{X}} f (x) d m (x), \end{array}
\]

从而 \({}^{t}h(m)=\lambda m\) 。这就证明了引理 4。

把 III, p. 78 的命题 4 与定理 4 应用于自同态 h，我们便得到下列命题。

定理 5. --- 设 \(\lambda\in K\) 。

\begin{enumerate}
\def\labelenumi{\alph{enumi})}
\item
  集 \(F_{\lambda}\) 与 \(G_{\lambda}\) 是 \(\mathcal{C}(X;K)\) 的有限维向量子空间，且它们的维数相等；
\item
  为使方程 (3) 至少有一个解 \(f \in \mathcal{C}(\mathrm{X}, \mathrm{K})\) ，必须且只需有 \(\int_{\mathrm{X}} a(x) g(x) d\mu(x) = 0\) （对每个函数 \(g \in G_{\lambda}\) ）。此时 (3) 的解集是 \(\mathcal{C}(\mathrm{X}; \mathrm{K})\) 的一个以 \(F_{\lambda}\) 为方向的仿射子空间；
\item
  下列互相排斥的条件之一成立：
\end{enumerate}

\begin{enumerate}
\def\labelenumi{(\roman{enumi})}
\item
  对任何函数 \(a \in \mathcal{C}(\mathrm{X};\mathrm{K})\) ，积分方程 (3) 存在唯一的解 \(f \in \mathcal{C}(\mathrm{X};\mathrm{K})\) ；
\item
  齐次方程 (4) 有一个非零解 \(f \in \mathcal{C}(\mathrm{X};\mathrm{K})\) 。
\end{enumerate}

推论. --- 由 \(\lambda\in\mathrm{K}\) （使空间 \(\mathrm{F}_{\lambda}\) 非零的）组成的集是 K 的一个可数、闭且离散的子集。

